% FORMALIZACION_NUEVA of the justification for the limit used in
% 70_radiacion_termica.tex. It does not introduce a constant generator.
\subsection{Thermodynamic limit of occupation sums}
\label{v:subsec:limite-termodinamico}

The preceding radial density can be obtained as an effective limit of the
sums in the periodic realization. The estimate must retain the fixed
constant mode and uniformly control both the origin and the spectral
tail. The action and calendar representations remain prior
to this realization operation.

\begin{lemma}[Convergence of the three thermal observables]
\label{v:lem:limite-termodinamico}
Let \(a,b>0\), \(L>0\), \(\Delta=2\pi/L\), and
\[
 F_N(r)=\frac{1}{e^{ar}-1},\qquad
 F_E(r)=\frac{br}{e^{ar}-1},\qquad
 F_Z(r)=-\log(1-e^{-ar})\quad(r>0).
\]
For each of these functions,
\begin{equation}
 \lim_{L\to\infty}\frac{2}{L^3}
 \sum_{\nu\in\mathbb Z^3\setminus\{0\}}
 F\!\left(\frac{2\pi}{L}|\nu|\right)
 =\frac{1}{\pi^2}\int_0^\infty r^2F(r)\,dr.
 \label{v:eq:limite-modos-discretos}
\end{equation}
The sums and the integral are finite. The specialization
\(a=\beta\hbar c\), \(b=\hbar c\) gives, respectively,
the number density, excitation-energy density, and logarithmic partition-function
density of the transverse modes.
\end{lemma}

\begin{proof}
All three functions are continuous and positive on \((0,\infty)\).
There exist constants \(C_0,C_1,d>0\), depending on \(a,b\) but
independent of \(\Delta\), such that
\begin{equation}
 F(r)\leq\frac{C_0}{r}\quad(0<r\leq1),\qquad
 F(r)\leq C_1e^{-dr}\quad(r\geq1).
 \label{v:eq:envolventes-termicas}
\end{equation}
Indeed, \(e^{ar}-1\geq ar\) bounds \(F_N\) and \(F_E\)
near the origin; moreover,
\(F_Z(r)=\log(1+F_N(r))\leq F_N(r)\leq1/(ar)\).
For \(r\geq1\), the inequality
\(-\log(1-y)\leq y/(1-y)\), with \(y=e^{-ar}\), gives the
exponential bound for \(F_Z\). The other two follow from
\((e^{ar}-1)^{-1}\leq e^{-ar}/(1-e^{-a})\), absorbing the
factor \(r\) in \(F_E\) into an exponential of smaller decay rate. These bounds
already prove integrability of \(r^2F(r)\).

For the discrete sums, group the lattice points by
\(m=\|\nu\|_\infty\geq1\). Each shell contains exactly
\[
 (2m+1)^3-(2m-1)^3=24m^2+2
\]
points and satisfies \(m\leq|\nu|\leq\sqrt3m\).
Fix \(0<\delta\leq1\) and \(0<\Delta\leq1\).
If \(\delta<\Delta\), there are no points with
\(0<\Delta|\nu|\leq\delta\). Otherwise, setting
\(N=\lfloor\delta/\Delta\rfloor\), the first bound in
\eqref{v:eq:envolventes-termicas} gives
\begin{align*}
 \Delta^3\!\sum_{0<\Delta|\nu|\leq\delta}F(\Delta|\nu|)
 &\leq C_0\Delta^2\sum_{m=1}^{N}\frac{24m^2+2}{m}\\
 &\leq26C_0\Delta^2\sum_{m=1}^{N}m
 \leq26C_0\delta^2.
\end{align*}
The second line uses \(m^{-1}\leq m\) and then
\(N(N+1)/2\leq N^2\) for \(N\geq1\).
The integral contribution from the same neighborhood is also of order
\(\delta^2\), since
\(\int_0^\delta r^2F(r)\,dr\leq C_0\delta^2/2\).

For \(R\geq1\), every point with \(\Delta|\nu|>R\) belongs
to a shell \(m>R/(\sqrt3\Delta)\). The second bound in
\eqref{v:eq:envolventes-termicas} gives
\begin{align*}
 \Delta^3\!\sum_{\Delta|\nu|>R}F(\Delta|\nu|)
 &\leq C_1e^{-dR/(2\sqrt3)}
 \Delta^3\sum_{m\geq1}(24m^2+2)e^{-d\Delta m/2}\\
 &\leq C_2e^{-dR/(2\sqrt3)},
\end{align*}
where \(C_2\) is independent of \(0<\Delta\leq1\). To verify
this uniformity, set \(q=e^{-d\Delta/2}\) and use
\[
 \sum_{m\geq1}q^m=\frac{q}{1-q},\qquad
 \sum_{m\geq1}m^2q^m=\frac{q(1+q)}{(1-q)^3}.
\]
The first identity is the geometric series, and the second follows by
applying \(q\,d/dq\) twice, which is valid for \(|q|<1\).
Moreover,
\(1-e^{-d\Delta/2}\geq\Delta(1-e^{-d/2})\) by concavity.
The factor \(\Delta^3\) therefore compensates for the denominators of
both expressions. The integral tail likewise tends to zero by the
exponential bound. For each \(\Delta>0\), the same estimates
prove convergence of the infinite sum.

On the compact annulus \(\delta\leq|k|\leq R\), the function
\(F(|k|)\) is continuous and bounded. Riemann sums of mesh
\(\Delta\) converge to its integral; the two boundary spheres
have measure zero. First let \(\Delta\to0\), then
\(\delta\to0\), and finally \(R\to\infty\), using the preceding
uniform bounds. The result is
\[
 \Delta^3\sum_{\nu\ne0}F(\Delta|\nu|)
 \longrightarrow \int_{\mathbb R^3}F(|k|)\,d^3k
 =4\pi\int_0^\infty r^2F(r)\,dr.
\]
Since \(2/L^3=2\Delta^3/(2\pi)^3\), the final radial coefficient
is \(2\cdot4\pi/(2\pi)^3=1/\pi^2\), proving
\eqref{v:eq:limite-modos-discretos}.
\end{proof}

Control at the origin is carried out over positive-energy excitations.
The constant mode remains fixed before the partition function is formed:
removing it afterward from an integral would not replace that condition on
the finite representation. The proof establishes the limit of the
specified realization and preserves the distinction between that realization and
the prior construction of its action and calendar parameters.
